\subsection{The normalized source certificate now has a Rocq proof}

The incoming normalization proposal already reduced the source-correspondence
certificate from $2240$ entries to $1426$, and supplied a Lean proof. The incoming
compact-source proposal already supplied the lossless positive-integer path codec
and a plain mirror theorem for Rocq. Our contribution combines those two lines of
work and fills an explicitly outstanding step: a kernel-checked Rocq theorem for
the \emph{same normalized certificate}, using the exact current source programs.
Neither the normalization idea nor its Lean implementation is claimed as new.

The resulting theorem is
\[
 \forall R,\qquad
 \operatorname{Halts}(p_{215},(\epsilon,R))
 \quad\Longleftrightarrow\quad
 \operatorname{Halts}(\mathrm{dagPerm},(\epsilon,R)),
\]
where $R$ is an arbitrary natural-valued register file. Both program definitions
are byte-for-byte equal to the definitions at the pinned repository revision.
The generic checker and the concrete source theorem print
\emph{Closed under the global context}. No consistency assumption occurs.
The tested adapter exports the existing qualified name
\texttt{K215F.Source215.p215\_dagPerm}; an explicit typed client checks that public
name and its original statement.

\paragraph{The mathematical reduction.}
For finite Boolean paths, put $P\sim Q$ when their infinite extensions by false
bits agree at every index. This is exactly the existing \texttt{ipeq} relation.
The existing \texttt{hn\_ipeq} theorem preserves halting after a specified number
of steps under this equivalence. The new Boolean comparator \texttt{ibeq} checks
$\sim$: after one path ends, every remaining bit of the other must be false.
Its soundness proof is an induction on the paths.

At a stored pair of decrement or halt sites, the normalized checker uses the
unchanged \texttt{seg} interpreter. Suppose its two successful segments have
increment lists $L_1,L_2$ and actual endpoint paths $Q_1,Q_2$. The checker verifies
\[
 \forall r,\quad \operatorname{count}(r,L_1)=\operatorname{count}(r,L_2),
 \qquad Q_1\sim P_1,\quad Q_2\sim P_2,
\]
where $(P_1,P_2)$ is the selected successor entry. The inherited segment lemmas
produce the actual source and target runs and the pointwise register equations
$R_i'(r)=R_i(r)+\operatorname{count}(r,L_i)$. Thus equal registers remain equal.
The source endpoint's halting run transfers to $P_1$ by \texttt{hn\_ipeq}; the
induction hypothesis gives a target run from $P_2$; the same equivalence theorem
transfers that run back to $Q_2$. Prepending the target segment completes the
branch argument.

The induction strictly decreases when the source performs its next decrement
test, irrespective of whether either connecting segment has length zero.
Both test outcomes are covered by the certificate. A halt entry terminates
immediately. Mirroring exchanges the two programs and the two pointers while
preserving both successor indices, so a generic theorem supplies the reverse
direction. Only one concrete checker evaluation is required.

\paragraph{What normalization does not authorize.}
The equivalence identifies only a finite path's \emph{trailing false bits} in
the existing mathematical pointer representation. It does not remove arbitrary
suffixes, change a consumed leaf path, or delete physical cells from a binary
Turing-machine tape. Break and continue still act on the exact consumed path
returned by the unchanged interpreter. This distinction is essential: a proof
about abstract pointer representations does not itself establish a different
tape-frame layout.

\paragraph{Size and data binding.}
The normalized certificate has $1426$ entries: $1425$ decrement checkpoints and
one halt checkpoint. There are $2851$ branch obligations including the initial
branch. Removing $814$ entries is a $36.3393\%$ reduction. The data plus the
retained codec occupy $109958$ uncompressed source bytes, compared with $173003$
for the incoming compact $2240$-entry presentation and $1499170$ for the original
Rocq offset certificate. The normalized data alone occupy $108633$ bytes; the
codec occupies $1325$ bytes. The new generic checker occupies $193$ lines and
$8315$ bytes, and the concrete source theorem occupies $21$ lines and $757$ bytes.
These figures measure source representation, not total proof-object size or
verification speed.

The two proof assistants share the exact JSON certificate with SHA-256
\begin{center}\small\ttfamily
6d30d89253629525b0af3066c7fec67b052c07413938fc7a94dd96ab1ff406e2
\end{center}
The incoming JSON writes pointer bits most-significant-first; the retained Rocq
codec recursively stores the head bit in the least-significant constructor.
The renderer performs this explicit conversion. An independent binding checker
compares all decoded paths and both successors in every row. This audit is
separate from the Rocq source theorem: even an incorrect producer cannot make
an invalid source correspondence pass the proved checker.
The positive-integer codec has no fixed decoding-fuel limit; the longest
$176$-bit pointer is represented without truncation.

\paragraph{Verification boundary and integration.}
The focused build compiles all twelve project modules needed for the new source
theorem, then the public adapter and its typed client. Its receipts record
commands, exact source hashes, raw output hashes, per-module times, and peak
resident memory. The toolchain is Coq $8.18.0$; the inherited, proved
\texttt{Compat818} module adapts standard-library lemma names in the private
copy of \texttt{RMEquiv}. No axiom or executable-definition change is introduced.
The original Rocq $9.1$ dependency source is retained unchanged for that
integration route. All twelve source modules compiled in $80.632$ seconds of summed compiler
wall time, with peak memory $824648$ KiB. A separate recursive \texttt{coqchk}
of the concrete source theorem completed successfully in $521.235$ seconds,
with exit code zero. The adapter and typed client took a further $4.478$ seconds.
The source module compiled in $13.718$ seconds in the recorded
run; its forward VM reduction reported $6.874$ seconds. These are observations,
not a controlled speed comparison with a baseline.

The new source theorem is a replacement component of the established machine
proof. It does not by itself constitute a fresh rebuild of the complete
$237$-module ZFC endpoint. The final integration step is to retain the public
source facade, update its generator and dependency manifests, rebuild
\texttt{Correct215} and \texttt{CorrectDense215}, and inspect their endpoint
assumption reports. The binary table and the logical calibration remain the
ones already present in the repository.
